\documentclass[12pt]{article}
\usepackage{amsmath,amssymb,amsfonts}
\usepackage[margin=1in]{geometry}

\begin{document}

\section*{Chapter 3, Problem 6}

\textbf{Problem.} Prove that a basis is a maximal linearly independent set of vectors.

\bigskip
\textbf{Solution.}

Let $V$ be an $n$-dimensional vector space and let $\mathcal{B} = \{x_1, x_2, \ldots, x_n\}$ be a basis for $V$. We must show that $\mathcal{B}$ is a \emph{maximal} linearly independent set, meaning:
\begin{enumerate}
    \item $\mathcal{B}$ is linearly independent.
    \item Adding any vector to $\mathcal{B}$ destroys linear independence.
\end{enumerate}

\medskip
\textbf{Part 1:} By definition, a basis is linearly independent.

\medskip
\textbf{Part 2:} Let $v \in V$ be any vector not in $\mathcal{B}$. Since $\mathcal{B}$ spans $V$, we can write
\[
v = \alpha_1 x_1 + \alpha_2 x_2 + \cdots + \alpha_n x_n
\]
for some scalars $\alpha_1, \ldots, \alpha_n$. This means
\[
1 \cdot v - \alpha_1 x_1 - \alpha_2 x_2 - \cdots - \alpha_n x_n = 0,
\]
which is a nontrivial linear relation among $\{v, x_1, x_2, \ldots, x_n\}$ (the coefficient of $v$ is 1, which is nonzero). Therefore the set $\mathcal{B} \cup \{v\}$ is linearly dependent.

Since no vector can be added to $\mathcal{B}$ while maintaining linear independence, $\mathcal{B}$ is a maximal linearly independent set. \hfill $\blacksquare$

\end{document}
